1. Abstract（要約：背景、問題、提案、定量的結果の要約を1–2文で）
2. Introduction（背景、課題、提案の概観、明確な貢献箇条書き）
3. Related Work（ACT‑R 他、トポロジー／ポテンシャル場、量子認知・量子インスパイア、カテゴリ理論の適用）
4. System Architecture（高レベル図）
- 4.1 Dual-Layer Overview (Control Hub vs Topological Field)
- 4.2 Sensory Processing & Particle Encapsulation
- 4.3 Topological Mapping & Spatial Allocation
- 4.4 Quantization & Entanglement Phase（定義・仮定）
5. Mathematical Formulation（ここを丁寧に）
- 5.1 記法と用語表（Notation table）
- 5.2 Functor / Category mapping（命題・補題があれば要約と証明の場所）
- 5.3 Potential Energy & Dynamics（微分方程式、収束条件）
- 5.4 Entanglement model（もし物理的意味があるなら厳密化）
6. Reasoning & Algorithms（実際の推論メカニズム）
- 6.1 Extended Match & Select（擬似コード/アルゴリズム）
- 6.2 Backtracking / Topological deformation（数式＋アルゴリズム）
7. Implementation & Complexity（再現性のために必須）
- 実装言語、ライブラリ、ハードウェア、計算複雑度、最適化手法
8. Experiments / Evaluation（新規で必ず追加）
- ベンチマーク（タスク例：推論精度、問題解決、学習速度、ロバスト性）
- 比較対象（ACT‑R, 既存のポテンシャル場法, neural baselines）
- 指標（accuracy、convergence time、memory、ablation）
- 結果（表＋図）、考察
9. Discussion（利点・限界・倫理的配慮・応用可能性）
10. Conclusion（まとめ・主な知見・将来の道筋）
Appendix（証明、ハイパーパラメータ、追加図表、ソースコード／データへのリンク）
